\ifdefined\XeTeXrevision\else\ifdefined\pdfoutput\pdfoutput=1\fi\fi
\documentclass[11pt]{article}
\ifdefined\XeTeXrevision\PassOptionsToPackage{xetex}{hyperref}\fi

\usepackage{amsmath,amssymb}
\usepackage{graphicx}
\usepackage[margin=1.1in]{geometry}
\usepackage[colorlinks=true,linkcolor=blue,citecolor=blue,urlcolor=blue]{hyperref}

\newcommand{\Keff}{K_{\mathrm{eff}}}
\newcommand{\bmin}{\beta_{\min}}
\newcommand{\dw}{\Delta\omega}
\newcommand{\taurelax}{\tau_{\mathrm{relax}}}

\title{Ghost Sensitivity Peaks in Adaptive Phase-Coupled Systems:\\
A Disguised Adler Bifurcation and a Calibration Rule for Coupling Floors}
\author{Mihai Ro\c{s}ca\\
\small Independent researcher, Br\u{a}ila, Romania\\
\small ORCID: \href{https://orcid.org/0009-0001-1422-6209}{0009-0001-1422-6209} \quad
\small \texttt{mihairosca1982@gmail.com}}
\date{July 2026 --- draft v0.1}

\begin{document}
\maketitle

\begin{abstract}
We study a Kuramoto system of $N$ internal oscillators coupled to an external
reference (``human intent'') through a noisy estimate and an adaptive coupling
weight $\beta$ governed by a coherence index (RSI). During calibration of a
coupling floor $\bmin$ --- introduced to prevent the adaptive loop from
self-isolating at high internal coherence --- the system's measured sensitivity
to genuine reference changes exhibited an unexpected non-monotonic peak at
$\bmin \approx 0.15$, reproduced across seven noise levels. Two intuitive
explanations (forced harmonic resonance; critical damping of the feedback loop)
were tested directly and rejected, one after identifying a fabricated data
column in an externally proposed analysis. The true mechanism is a
saddle-node-on-invariant-circle (SNIC) phase-locking bifurcation of Adler type
in the mean-field reduction: near the locking threshold $\Keff = \dw$, critical
slowing down makes fixed-window transient measurements register their deepest
dips, producing a \emph{ghost peak} that tracks the observation protocol rather
than any resonance of the dynamics. The scaling law
$\taurelax = (\Keff^2 - \dw^2)^{-1/2}$ is confirmed on the isolated Adler
equation (measured exponents $-0.4957$ and $-0.5009$ in two independent
implementations) and on the full 30-oscillator system (measured-to-predicted
ratios 0.985--0.999). All quantitative claims were reproduced by an independent
reimplementation with different seeds, RNG, and integration scheme. We derive a
calibration rule for coupling floors, $\bmin \ge m\,\dw_{\max}/K_{\mathrm{ext}}$
with $m(\tau_{\mathrm{target}}) = \sqrt{1 + (\dw_{\max}\tau_{\mathrm{target}})^{-2}}$,
and discuss a general methodological warning: benchmarks of adaptive alignment
mechanisms that use fixed observation windows can manufacture spurious optima
near bifurcations.
\end{abstract}

\section{Introduction}

Adaptive coupling schemes --- systems that adjust how strongly they follow an
external reference based on their own estimate of coherence --- arise in AI
alignment proposals \cite{miyato2025}, robot swarms, and power-grid control. A
recurring design element is a \emph{coupling floor}: a minimum weight below
which the follow-the-reference term is never allowed to fall, guarding against
a failure mode in which a self-confident system stops listening precisely
because it is coherent.

This paper documents a case study in calibrating such a floor. The phenomenon
that motivated it looked, at first, like a discovery: sensitivity to genuine
reference changes peaked at an intermediate floor value instead of increasing
monotonically, consistently across noise levels. We show that the peak is not
a resonance, not a critical-damping effect, and not a property of the adaptive
rule at all. It is the shadow of a textbook SNIC bifurcation
\cite{adler1946,strogatzbook,ermentrout1986} cast onto a fixed observation
window by critical slowing down --- a \emph{ghost peak}.

Contributions:
\begin{enumerate}
\item \textbf{Mechanism identification.} The mean-field reduction of an
  adaptively coupled Kuramoto system with a coupling floor is an Adler
  equation; the empirical sensitivity peak is explained completely by critical
  slowing down near its locking threshold (Section~\ref{sec:results}).
\item \textbf{Quantitative verification at two levels.} The
  $\tau \sim (\Keff - \dw)^{-1/2}$ divergence is confirmed on the isolated
  Adler equation and on the full system, with ratios 0.985--0.999
  (Fig.~\ref{fig:scaling}), including an independent reimplementation with
  different seeds and integrator.
\item \textbf{A calibration rule with a derived margin.}
  $\bmin \ge m\,\dw_{\max}/K_{\mathrm{ext}}$, where the safety margin $m$
  follows from a recovery-time requirement:
  $m(\tau_{\mathrm{target}}) = \sqrt{1+(\dw_{\max}\tau_{\mathrm{target}})^{-2}}$
  (Section~\ref{sec:calibration}).
\item \textbf{A methodological warning.} Fixed-window sensitivity benchmarks
  near bifurcations produce spurious optima. Two plausible-sounding
  alternative explanations failed direct testing; one arrived with an
  algebraically derived (not measured) data column, which the verification
  protocol caught (Section~\ref{sec:rejected}).
\end{enumerate}

We are explicit about what this work is not: it is a nonlinear-dynamics case
study on a small model ($N=30$) motivated by AI-alignment questions, not
evidence about production-scale alignment methods (Section~\ref{sec:limitations}).

\section{Model}\label{sec:model}

$N$ internal phase oscillators $\theta_i$ with natural frequencies $\omega_i$
couple internally (Kuramoto) and to an estimate of an external reference phase
$\theta_h$:
\begin{equation}\label{eq:model}
\frac{d\theta_i}{dt} = \omega_i
 + \alpha\,\frac{K_{\mathrm{int}}}{N}\sum_j \sin(\theta_j-\theta_i)
 + \beta\,K_{\mathrm{ext}}\sin(\hat\theta_h-\theta_i).
\end{equation}
The reference is observed only through noisy proxies
$z(t) = \theta_h(t) + \xi(t)$, $\xi \sim \mathcal{N}(0,\sigma^2)$, filtered by
a constant-velocity Kalman filter to produce $\hat\theta_h$. Coherence is
monitored by
\begin{equation}\label{eq:coherence}
\Phi_{\mathrm{int}} = \bigl|\langle e^{i\theta_j}\rangle\bigr|, \qquad
\Phi_{\mathrm{ext}} = \tfrac{1}{2}\bigl(1+\cos(\bar\theta-\hat\theta_h)\bigr), \qquad
\Phi = \alpha\,\Phi_{\mathrm{int}} + \beta\,\Phi_{\mathrm{ext}},
\end{equation}
where $\bar\theta = \arg\langle e^{i\theta_j}\rangle$. (An earlier draft of the
framework used $\Phi_{\mathrm{ext}} = |e^{i(\bar\theta-\hat\theta_h)}|$, which
is identically 1; the degeneracy was found by running the simulation, not by
inspection.) The Resonance Stability Index RSI is a 50-step moving average of
$\Phi$, and the adaptive weights are
\begin{equation}\label{eq:adaptive}
\beta = \max(1-\mathrm{RSI},\,\bmin), \qquad \alpha = 1-\beta.
\end{equation}
Equation~\eqref{eq:adaptive} implements the design intent ``open toward the
reference when unstable''; $\bmin$ is the coupling floor whose calibration this
paper concerns. Figure~\ref{fig:system} shows the architecture.

\begin{figure}[t]
\centering
\includegraphics[width=0.78\linewidth]{fig1_system.pdf}
\caption{Architecture. $N$ internal Kuramoto oscillators; an external reference
observed through noisy proxies and a Kalman filter; adaptive weights
$\alpha,\beta$ from the RSI loop with floor $\bmin$.}
\label{fig:system}
\end{figure}

\paragraph{Mean-field reduction.} When the internal population is strongly
synchronized ($\Phi_{\mathrm{int}} \approx 1$), the group behaves as one
effective oscillator, and the phase difference $\psi = \bar\theta - \theta_h$
obeys the Adler equation \cite{adler1946}
\begin{equation}\label{eq:adler}
\frac{d\psi}{dt} = -\dw - \Keff\sin\psi, \qquad \Keff = \beta\,K_{\mathrm{ext}},
\end{equation}
with $\dw$ the internal--reference frequency mismatch. For $\Keff > |\dw|$
there is a stable lock $\psi^* = -\arcsin(\dw/\Keff)$ with relaxation time
\begin{equation}\label{eq:tau}
\taurelax = \bigl(\Keff^2 - \dw^2\bigr)^{-1/2};
\end{equation}
for $\Keff < |\dw|$ the phase slips with period
$T_{\mathrm{slip}} = 2\pi\,(\dw^2-\Keff^2)^{-1/2}$. Both diverge at
$\Keff = \dw$ --- a SNIC bifurcation with critical slowing down
\cite{strogatzbook,ermentrout1986} (Fig.~\ref{fig:bifurcation}).

\begin{figure}[t]
\centering
\includegraphics[width=0.62\linewidth]{fig2_bifurcation.pdf}
\caption{Adler locking diagram. Stable lock $\psi^* = -\arcsin(\dw/\Keff)$
(solid) and unstable branch (dashed) merge in a SNIC bifurcation at
$\Keff = \dw$; no fixed point exists in the shaded slipping region.}
\label{fig:bifurcation}
\end{figure}

\section{Methods}\label{sec:methods}

\paragraph{Simulation protocol (calibration experiment).} $N=30$,
$\omega_i \sim \mathcal{N}(1.0,\,0.05)$, $dt=0.02$, $T=3000$ steps,
$K_{\mathrm{int}}=1.2$, $K_{\mathrm{ext}}=1.5$, Euler integration. The
reference frequency $\omega_h$ switches $1.0 \to 1.3$ at $t=15$
($\dw_1 = 0.3$) and $1.3 \to 0.9$ at $t=35$ ($\dw_2 = 0.4$). Kalman:
$Q = \mathrm{diag}(10^{-4},10^{-3})$, $R = \max(\sigma^2, 10^{-4})$. Grid:
$\bmin \in \{0, 0.02, 0.05, 0.08, 0.10, 0.15, 0.20, 0.30\}$ $\times$
$\sigma \in \{0.1, 0.25, 0.5, 1.0, 1.5, 2.0, 2.5\}$, seeds $\{1,2,3\}$.

\paragraph{Sensitivity metric (formal definition).} With switch times
$t_1, t_2$ and a fixed window $T_w = 250$ steps (5\,s):
\begin{equation}\label{eq:metric}
S(\bmin,\sigma) = \frac{1}{2}\sum_{k=1}^{2}
\Bigl[\,1 - \min_{t\in[t_k,\,t_k+T_w]} \Phi_{\mathrm{ext}}(t)\Bigr].
\end{equation}
$S$ is a \emph{transient} quantity by construction; steady-state
$\Phi_{\mathrm{ext}}$ returns to $\approx 0.99$ in all locked regimes.
\textbf{Jitter} is the standard deviation of $d\bar\theta/dt$ over a
pre-switch stationary segment (steps 100--700), quantifying noise transmission
to the group phase.

\paragraph{Scaling measurements.} $\taurelax$ is measured (i) on
Eq.~\eqref{eq:adler} directly (RK4, $dt=10^{-3}$), from the asymptotic slope of
$\ln|\psi-\psi^*|$, for $\Keff/\dw - 1 \in [10^{-3},10^{-1}]$; and (ii) on the
full system (RK4, $dt=5\times10^{-3}$, fixed $\beta$, spread
$\sigma_\omega = 0.01$) by applying a 0.2\,rad phase kick to the reference
after settling and fitting the exponential return of $\psi$. Two measurement
bugs found and fixed during this step are documented in
Appendix~\ref{app:record} (phase-wrapping of $\psi$; sign convention of $\dw$)
--- before the fixes, the full-system check appeared to fail by a factor 4--14.

\paragraph{Independent reproduction.} All headline numbers were re-derived by
a second implementation written without access to the original code
(vectorized mean-field formulation instead of pairwise sums, independently
written Kalman filter, a different RNG, seeds $\{10,11,12\}$). Where the two
implementations disagree is stated explicitly in Results.

\section{Results}\label{sec:results}

\subsection{The phenomenon: a non-monotonic sensitivity peak}

$S(\bmin)$ rises from $\approx 0.29$ at $\bmin=0$ to a maximum of 0.57--0.68 at
$\bmin=0.15$, then falls to $\approx 0.12$--$0.36$ by $\bmin=0.30$,
consistently across all seven noise levels and in both implementations
(Fig.~\ref{fig:ghost}). Naively, more guaranteed coupling should mean more
sensitivity; the peak demands a mechanism.

\subsection{Two plausible explanations, tested and rejected}\label{sec:rejected}

\paragraph{Forced harmonic resonance.} Hypothesis: the RSI averaging window
(50 steps) resonates with the switching cadence. Test: periodic forcing at
four periods (10--80 steps). Result: response amplitude decreases
monotonically with $\bmin$ at every period --- no peak. Rejected.

\paragraph{Critical damping of the RSI--$\beta$ loop.} Hypothesis: the peak
marks a critically damped regime of the feedback loop. The supporting table
supplied with this hypothesis contained a ``steady-state response'' column
that direct measurement exposed as algebraically derived from the transient
column ($\approx 1 - \text{transient}$) rather than measured: actual
steady-state $\Phi_{\mathrm{ext}}$ is $\approx 0.99$ for \emph{all} $\bmin$.
Rejected; the episode is retained here as part of the verification record.

\subsection{The mechanism: a disguised Adler bifurcation}

The two switching events have thresholds
$\beta_{\mathrm{thr},1} = \dw_1/K_{\mathrm{ext}} = 0.20$ and
$\beta_{\mathrm{thr},2} = \dw_2/K_{\mathrm{ext}} \approx 0.267$
(Eq.~\eqref{eq:adler}). Below both thresholds the system slips persistently;
far above both it relocks so fast that a 5\,s window barely registers a dip.
Near threshold, critical slowing down (Eq.~\eqref{eq:tau}) makes transients
longer than the window --- which therefore records its deepest minima
\emph{near but not at} threshold. The measured peak ($\bmin \approx 0.12$--$0.15$,
just below $\beta_{\mathrm{thr},1}$, with the RSI--$\beta$ loop supplying
transient coupling above the floor) is the fixed window's shadow of the
bifurcation, not a property of the adaptive rule. The peak location is
protocol-dependent by construction --- hence \emph{ghost peak}.

\begin{figure}[t]
\centering
\includegraphics[width=0.7\linewidth]{fig4_ghost_peak.pdf}
\caption{The ghost peak. Transient sensitivity (Eq.~\eqref{eq:metric}) vs
$\bmin$ at $\sigma=0.25$: original implementation (seeds 1--3) and independent
reimplementation (seeds 10--12; error bars $= \pm 1$ std over seeds). Dashed
lines: Adler thresholds of the two switching events. The peak sits below
threshold, in the critical-slowing region --- its location tracks the
measurement protocol, not a resonance.}
\label{fig:ghost}
\end{figure}

\subsection{Quantitative verification}

\begin{itemize}
\item \textbf{Locking threshold (full system):} persistent slipping below,
  definitive locking above; empirical transition at $\bmin \approx 0.20$--$0.25$
  vs analytic 0.20, the offset explained by the RSI--$\beta$ loop raising
  $\beta$ above the floor when coherence drops.
\item \textbf{Scaling exponent (isolated Adler):} $-0.5009$ (original) and
  $-0.4957$ (independent reimplementation) vs $-0.5$ analytic;
  $T_{\mathrm{slip}}$ measured/predicted $= 1.0000$ across
  $\Keff/\dw \in \{0.5, 0.8, 0.95, 0.99\}$.
\item \textbf{Transfer to the full system ($N=30$):} measured/predicted $\tau$
  ratios 0.9847--0.9965 (original, larger $\omega$-spread) and 0.9968--0.9990
  (reimplementation, $\sigma_\omega = 0.01$) across
  $\beta \in \{0.35, 0.4, 0.5, 0.7, 1.0\}$ (Fig.~\ref{fig:scaling}). The
  residual deviation shrinks with the internal frequency spread, consistent
  with the mean-field idealization.
\item \textbf{Ghost peak reproduction:} the $\bmin = 0.15$ peak reproduces at
  all seven noise levels under the independent reimplementation (peak
  $S = 0.55$--$0.67$ vs neighbors 0.25--0.43), with different seeds and RNG
  (Fig.~\ref{fig:ghost}).
\end{itemize}

\begin{figure}[t]
\centering
\includegraphics[width=0.72\linewidth]{fig3_scaling.pdf}
\caption{Critical slowing down. Relaxation time vs distance from threshold,
log-log: isolated Adler equation (circles), full $N=30$ system (squares),
analytic curve (line), $-1/2$ slope guide (dotted). Both systems collapse onto
Eq.~\eqref{eq:tau}.}
\label{fig:scaling}
\end{figure}

\subsection{Calibration rule with a derived margin}\label{sec:calibration}

Requiring the locked system to recover within a target time,
$\taurelax \le \tau_{\mathrm{target}}$, and writing $\Keff = m\,\dw_{\max}$,
Eq.~\eqref{eq:tau} gives
\begin{equation}\label{eq:margin}
m(\tau_{\mathrm{target}}) = \sqrt{1 + \bigl(\dw_{\max}\,\tau_{\mathrm{target}}\bigr)^{-2}},
\qquad\text{equivalently}\qquad
\taurelax \le \bigl[\dw_{\max}\sqrt{m^2-1}\bigr]^{-1}.
\end{equation}
The rule $\bmin \ge m\,\dw_{\max}/K_{\mathrm{ext}}$ with $m=1.5$ therefore
guarantees $\taurelax \le 0.894/\dw_{\max}$ --- recovery within one
characteristic mismatch time. A numerical check on Eq.~\eqref{eq:adler} gives
guaranteed vs measured $\tau$ agreement of 0.996--0.999 for
$m \in \{1.2, 1.5, 2, 3\}$. This replaces both the naive ``$\bmin$ that scored
best on the grid'' (0.15 --- inside the critical-slowing region, hence
fragile) and any fixed universal constant: the floor must be set from the
anticipated frequency mismatch and the required recovery time, quantities an
implementer can estimate.

\section{Discussion}\label{sec:discussion}

\paragraph{For benchmark design.} Any evaluation of an adaptive
coupling/alignment mechanism that (i) measures transient responses (ii) in
fixed windows (iii) while sweeping a parameter that moves the system across a
locking threshold will manufacture an interior optimum. The optimum tracks the
observation protocol --- window length, event magnitudes $\dw_k$ --- not the
mechanism. This is the transient counterpart of using critical slowing down as
an early-warning signal \cite{scheffer2009}: the same divergence that provides
warning signals also corrupts fixed-window benchmarks.

\paragraph{For adaptive alignment schemes.} The episode illustrates a
structural risk in coherence-gated coupling (Eq.~\eqref{eq:adaptive}): the
regime in which the system most needs external correction (large $\dw$) is
exactly where slow relaxation makes corrections hardest to observe. Recovery
time after a large perturbation deserves treatment as a first-class metric,
separate from steady-state alignment quality.

\paragraph{On the verification record.} Two attractive explanations failed
direct tests; one external contribution contained fabricated data; two bugs in
our own verification code initially masqueraded as a 4--14$\times$ violation
of the mean-field prediction. We retain these in Appendix~\ref{app:record}
because the paper's reliability claim rests on the process, not on the
elegance of the final formula.

\section{Limitations}\label{sec:limitations}

\begin{itemize}
\item Results characterize a small phase-oscillator model ($N=30$) with
  synthetic reference signals. No claim is made about neural-network training
  dynamics or production alignment methods (RLHF/DPO); the connection is
  motivational, not demonstrated.
\item The proxies-to-phase pipeline (Kalman with fixed $R$) is idealized;
  whether real human-interaction signals admit a useful phase representation
  is an open problem and the framework's main unvalidated assumption.
\item The ghost-peak location depends on protocol parameters ($T_w$, $\dw_k$);
  we characterize the mechanism, not a universal constant.
\item Euler integration with $dt=0.02$ for the calibration experiment (RK4 for
  scaling measurements); the peak's robustness to integrator choice is
  confirmed only indirectly through the independent reimplementation.
\end{itemize}

\section{Reproducibility and code availability}

Two independent implementations produce all headline numbers: the original
(pairwise-sum Kuramoto, legacy seeding, seeds 1--3) and a reimplementation
written from the equations alone (vectorized mean field, independent Kalman,
seeds 10--12). Scripts, exact parameters, and both raw outputs accompany this
manuscript: code repository
\url{https://github.com/amidigiart/p6-adler-ghost-peak}, archived at
\href{https://doi.org/10.5281/zenodo.21269201}{doi:10.5281/zenodo.21269201}.
The original calibration script reproduces its archived output bit-for-bit
under Python 3.14 / NumPy 2.5.

\section*{Acknowledgments}

The author thanks, in advance, the first readers and endorsers of this
manuscript: the willingness of established researchers to examine work
arriving from outside institutional walls is what makes independent research
possible, and it is not taken for granted. This work is dedicated to my son,
Patrick --- it was built so that anyone, at any time, can verify it without
having to trust its author.

\section*{AI-assisted research disclosure}

The research was conducted by the author in iterative sessions with multiple
AI systems used as drafting, critique, and verification instruments, with
executed code as the arbiter for every quantitative claim. Two AI-proposed
explanations were rejected by direct testing, and one AI-supplied table was
identified as containing fabricated values (Section~\ref{sec:rejected}). The
independent reimplementation and cross-validation were performed by an AI
system (Claude, Anthropic) in a separate session without access to the
original code. All results were verified by execution on the author's
hardware.

\appendix
\section{Verification record (bugs and rejected analyses)}\label{app:record}

\paragraph{A.1 Degenerate coherence formula.} The original external-coherence
definition $|e^{i(\bar\theta-\hat\theta_h)}| \equiv 1$ cannot detect
misalignment; found by simulation (RSI pinned at 0.999 across intent changes),
fixed to Eq.~\eqref{eq:coherence}.

\paragraph{A.2 Phase-wrapping bug.} $\psi$ computed as (wrapped $\bar\theta$)
$-$ (unwrapped $\theta_h$) inflated full-system $\tau$ ratios to
4.3--14.5$\times$; fixed by wrapping the difference.

\paragraph{A.3 Sign convention.} The implemented coupling direction yields
$d\psi/dt = -\dw - \Keff\sin\psi$; the initially assumed $+\dw$ made
predictions incomparable.

\paragraph{A.4 Fabricated data column.} An externally supplied ``empirical''
table for the critical-damping hypothesis contained a steady-state column
derived algebraically from the transient column; direct measurement
(steady-state $\Phi_{\mathrm{ext}} \approx 0.99\ \forall \bmin$) exposed it.

\paragraph{A.5 Rejected calibrator.} An externally supplied $\bmin$ calibrator
produced hardcoded recommendations that failed cross-validation against full
simulations (e.g., predicted $\Phi_{\mathrm{ext}} = 0.66$ vs measured 0.998 at
$\sigma = 0.05$).

\begin{thebibliography}{9}
\bibitem{adler1946} R.~Adler, ``A study of locking phenomena in oscillators,''
  \emph{Proc. IRE} \textbf{34}, 351--357 (1946).
\bibitem{kuramoto1984} Y.~Kuramoto, \emph{Chemical Oscillations, Waves, and
  Turbulence} (Springer, 1984).
\bibitem{strogatz2000} S.~H. Strogatz, ``From Kuramoto to Crawford: exploring
  the onset of synchronization in populations of coupled oscillators,''
  \emph{Physica D} \textbf{143}, 1--20 (2000).
\bibitem{strogatzbook} S.~H. Strogatz, \emph{Nonlinear Dynamics and Chaos}
  (Addison-Wesley, 1994), Ch.~4.
\bibitem{ermentrout1986} G.~B. Ermentrout and N.~Kopell, ``Parabolic bursting
  in an excitable system coupled with a slow oscillation,'' \emph{SIAM J.
  Appl. Math.} \textbf{46}, 233--253 (1986).
\bibitem{scheffer2009} M.~Scheffer \emph{et al.}, ``Early-warning signals for
  critical transitions,'' \emph{Nature} \textbf{461}, 53--59 (2009).
\bibitem{miyato2025} T.~Miyato, S.~L\"owe, A.~Geiger, and M.~Welling,
  ``Artificial Kuramoto Oscillatory Neurons,'' \emph{ICLR} (2025).
\bibitem{pikovsky2001} A.~Pikovsky, M.~Rosenblum, and J.~Kurths,
  \emph{Synchronization: A Universal Concept in Nonlinear Sciences}
  (Cambridge Univ. Press, 2001).
\end{thebibliography}

\end{document}
